% v2026a baseline; generated by R/44_Ideology_Diagnosis.R; SE clustered by party
\begin{table}[ht]
\caption{The populism $\times$ ideology interaction across five ideology measures}
\label{tab:oa2f}
\footnotesize
\centering
\resizebox{\ifdim\width>\linewidth\linewidth\else\width\fi}{!}{%
\begin{tabular}{lrrrrrl}
\toprule
Outcome & IDEO (LLM) & CMP rile & CMP rile (logit) & CHES lrgen & ParlGov left-right & Verdict \\
\midrule
Liberalism (overall) & $-$0.032* & $-$0.011 & $-$0.006 & $-$0.070* & $-$0.007 & no significant difference with any measure \\
Political & $-$0.060*** & $-$0.032* & $-$0.019 & $-$0.061 & 0.014 & right steeper in 1 of 5 measures \\
Social & $-$0.162*** & $-$0.076*** & $-$0.053* & $-$0.093* & $-$0.057** & right steeper in 3 of 5 measures \\
Economic & 0.075*** & 0.047** & 0.038* & $-$0.048 & 0.013 & right flatter in 2 of 5 measures \\
Financial markets & 0.027 & $-$0.003 & $-$0.015 & $-$0.103*** & $-$0.034 & right steeper in 1 of 5 measures \\
\midrule
N & 2,999--3,192 & 3,118--3,340 & 3,111--3,333 & 829--873 & 1,862--2,005 & \\
\bottomrule
\end{tabular}}
\par\smallskip
\begin{minipage}{\linewidth}\scriptsize Coefficient on populism $\times$ ideology in column (5) of Tables 4--8 (country and year fixed effects, SE clustered by party), with the ideology measure in the column standardised to mean 0 and SD 1. The samples differ by measure, because each external measure covers a different subset of manifestos (N). Verdict: number of measures with a significant ($p<0.05$) negative (right-leaning parties: steeper decline) or positive (flatter) interaction. *** p$<$0.01, ** p$<$0.05, * p$<$0.1.\end{minipage}
\end{table}
